% ============================================================
%  پیوست ب: تشریح تکمیلی — انسان‌شناسی، تنش‌های ساختاری
% ============================================================
\section{پیوست ب: تشریح تکمیلی نکات تحلیلی}
\label{app:analysis}

% ─────────────────────────────────────────────────────────────
\subsection{ب.۱ — انسان‌شناسی زیربنایی هر مکتب}
\label{subsec:anthropology}
% ─────────────────────────────────────────────────────────────

\begin{tcolorbox}[defbox, title={مقدمه: هر نظریه‌ی آزادی یک تصویر از انسان دارد}]
پیش‌فرض‌های انسان‌شناختی تعیین‌کننده‌ترین عامل در شکل‌گیری نظریه‌های آزادی هستند. اگر انسان را موجودی اجتماعی بدانی یا اتمیستیک، خودبسنده یا وابسته، عاقل یا احساساتی—پاسخ تو به «آزادی چیست» کاملاً تغییر خواهد کرد.
\end{tcolorbox}

\vspace{0.5cm}

\begin{table}[h!]
\centering
\caption{انسان‌شناسی زیربنایی مکاتب اصلی آزادی}
\label{tab:anthropology}
\renewcommand{\arraystretch}{1.8}
\setlength{\tabcolsep}{5pt}

\begin{tabularx}{\textwidth}{|>{\bfseries\color{PrimaryMid}}p{2.4cm}
  |>{\raggedright}X|>{\raggedright}X|>{\raggedright}X
  |>{\raggedright\arraybackslash}X|}

\hline
\rowcolor{PrimaryDeep}
\textcolor{white}{\textbf{مکتب}} &
\textcolor{white}{\textbf{تصویر انسان}} &
\textcolor{white}{\textbf{توانایی اصلی}} &
\textcolor{white}{\textbf{خطر اصلی برای انسان}} &
\textcolor{white}{\textbf{نتیجه‌ی سیاسی}} \\
\hline

\rowcolor{BgCool}
لیبرالیسم کلاسیک &
اتم عاقل و خودبسنده &
محاسبه‌ی منفعت شخصی &
مداخله‌ی دولت &
دولت حداقلی + بازار \\
\hline

\rowcolor{white}
لیبرالیسم ارزش‌مبنا &
موجود اجتماعی با پروژه‌ی خودشکوفایی &
اراده‌ی اخلاقی + عقل &
محرومیت + اجبار &
دولت رفاه محدود \\
\hline

\rowcolor{BgCool}
جمهوری‌خواهی &
شهروند سیاسی نیازمند مشارکت &
کنش در فضای عمومی &
سلطه‌ی خصوصی و عمومی &
نهادهای مدنی قوی \\
\hline

\rowcolor{white}
مارکسیسم &
موجود طبقاتی–تولیدی &
کار خلاق غیربیگانه &
استثمار اقتصادی &
انقلاب + جامعه‌ی بی‌طبقه \\
\hline

\rowcolor{BgCool}
فمینیسم &
موجود جنسیتی–بدنمند &
عاملیت زنانه &
پدرسالاری + خشونت &
سیاست جنسیتی \\
\hline

\rowcolor{white}
اگزیستانسیالیسم &
موجود رها در وجود + بدون ماهیت ثابت &
انتخاب اصیل &
فرار از مسئولیت (بدایمانی) &
اخلاق فردی \\
\hline

\rowcolor{BgCool}
رویکرد قابلیت &
موجود با قابلیت‌های خاص، متفاوت &
شکوفایی چندبُعدی &
محرومیت قابلیتی &
سیاست توسعه‌ی انسانی \\
\hline

\rowcolor{white}
سنت اسلامی (سنتی) &
عبد–خلیفه (بنده و جانشین خدا) &
تکلیف الهی &
گناه + طغیان &
شریعت + امارت \\
\hline

\rowcolor{BgCool}
سنت اسلامی (اصلاح‌طلب) &
موجود کرامتمند با اجتهاد &
عقل + ایمان &
استبداد + استعمار &
دموکراسی+اسلام \\
\hline

\rowcolor{white}
پسامدرنیسم (فوکو) &
موجود ساخته‌شده توسط گفتمان &
مقاومت و بازتعریف خود &
نرمال‌سازی قدرت &
سیاست‌های هویتی \\
\hline

\end{tabularx}
\end{table}

% ─────────────────────────────────────────────────────────────
\subsection{ب.۲ — تنش‌های ساختاری: پنج بحران حل‌نشده}
\label{subsec:tensions}
% ─────────────────────────────────────────────────────────────

\begin{tcolorbox}[enemybox, title={تنش ساختاری اول: آزادی فردی در برابر آزادی جمعی}]

\textbf{صورت‌بندی:} تا کجا آزادی فرد مجاز است بر جامعه تأثیر بگذارد؟ اصل آسیب میل می‌گوید: تنها زمانی که به دیگران آسیب می‌زند. اما «آسیب» چه معنایی دارد؟ آیا فساد اخلاقی جامعه آسیب است؟ آیا برخی انتخاب‌های سبک زندگی بر فرهنگ مشترک آسیب می‌زند؟

\textbf{گرهگاه:} این تنش در مباحث آزادی بیان (نفرت‌پراکنی)، آزادی مذهب، و آزادی سبک زندگی امروز کاملاً زنده است.

\textbf{پاسخ‌های رقیب:}
\begin{itemize}[leftmargin=2em]
  \item \textit{لیبرالیسم:} تنها آسیب‌های مستقیم مجاز هستند محدودیت را توجیه کنند.
  \item \textit{جامعه‌گرایی:} آزادی اجتماعی—نه فردی—اصل است؛ جامعه مجاز به محدود کردن خودفروپاشی است.
  \item \textit{پتیت:} تداخل تنها زمانی ناقض آزادی است که دلبخواهانه باشد—نه آنگاه که بر اساس قانون و با مشارکت همگان صورت گیرد.
\end{itemize}
\end{tcolorbox}

\begin{tcolorbox}[enemybox, title={تنش ساختاری دوم: آزادی صوری در برابر آزادی واقعی}]

\textbf{صورت‌بندی:} حق رأی یکسان برای همه وجود دارد. اما کسی که بی‌سواد است، در فقر زندگی می‌کند، و رسانه‌ها دسترسی ندارد—آیا واقعاً آزاد است که رأی بدهد؟

\textbf{گرهگاه:} تمایز «آزادی از» (رفع مانع قانونی) و «آزادی برای» (توانایی واقعی) اینجا نقطه‌ی تقابل اصلی است.

\textbf{پاسخ‌های رقیب:}
\begin{itemize}[leftmargin=2em]
  \item \textit{لیبرتاریانیسم:} آزادی صوری کافی است؛ نابرابری نتیجه‌ی انتخاب‌های آزاد است.
  \item \textit{رالز:} اولین اصل آزادی‌های اساسی را تضمین می‌کند، اما اصل دوم از نابرابری‌های مضر جلوگیری می‌کند.
  \item \textit{سن:} بدون قابلیت‌های واقعی، آزادی صوری بی‌معناست.
\end{itemize}
\end{tcolorbox}

\begin{tcolorbox}[enemybox, title={تنش ساختاری سوم: آزادی و اقتدار}]

\textbf{صورت‌بندی:} هر نظام سیاسی نوعی اقتدار دارد. حتی لیبرالیسم حداقلی به دولت نیاز دارد تا حقوق را اجرا کند. پس نظم اجتماعی همیشه «آزادی» را محدود می‌کند. سؤال این نیست که آیا محدودیت باشد، بلکه \textit{کدام محدودیت} و \textit{از طریق چه روندی}.

\textbf{گرهگاه:} مشروعیت اقتدار چه رابطه‌ای با آزادی دارد؟
\begin{itemize}[leftmargin=2em]
  \item \textit{هابز:} اقتدار قوی شرط لازم آزادی است—بدون لویاتان، آزادی طبیعی به جنگ همه علیه همه تبدیل می‌شود.
  \item \textit{لاک:} اقتدار تنها وقتی مشروع است که از رضایت مبتنی بر حق طبیعی برآمده باشد.
  \item \textit{آرنت:} اقتدار و آزادی هر دو به فضای عمومی نیاز دارند؛ یکی بدون دیگری ممکن نیست.
\end{itemize}
\end{tcolorbox}

\begin{tcolorbox}[enemybox, title={تنش ساختاری چهارم: آزادی و هویت}]

\textbf{صورت‌بندی:} آیا آزادی مستلزم «خود بدون تعهد» (Unencumbered Self) است—چنان‌که لیبرالیسم کانتی فرض می‌کند—یا ما همیشه درون یک سنت، فرهنگ، و هویت آزادیم؟

\thinker{مایکل سندل} استدلال کرد که «خود بدون تعهد» توهمی است: ما از پیش در جوامع اخلاقی فرو رفته‌ایم که آزادی ما را شکل می‌دهند.\footnote{\lr{Sandel, M. (1982). \textit{Liberalism and the Limits of Justice}. Cambridge University Press, p. 179.}} % FIXED: \lr in footnote

\textbf{پاسخ لیبرال:} ما می‌توانیم از هویت‌های خود انتقادی داشته باشیم. آزادی نه خلاء است، بلکه توانایی بازنگری در تعهدات است.
\end{tcolorbox}

\begin{tcolorbox}[enemybox, title={تنش ساختاری پنجم: آزادی و مسئولیت}]

\textbf{صورت‌بندی:} اگر آزادی واقعی است، پس مسئولیت کامل هم واقعی است. اما اگر شرایط اجتماعی، اقتصادی و فرهنگی انتخاب‌های ما را محدود می‌کنند، آیا «مسئولیت کامل» عادلانه است؟

\textbf{سارتر:} هیچ بهانه‌ای قابل قبول نیست. «محکوم به آزادی» یعنی محکوم به مسئولیت کامل.

\textbf{مارکسیسم:} این ایدئولوژی «مسئولیت شخصی» است که فقر ساختاری را توجیه می‌کند.

\textbf{رالز:} جایگاه اجتماعی-اقتصادی را انتخاب نکرده‌ایم؛ عدالت باید «شانس‌های دلبخواهانه» را خنثی کند.
\end{tcolorbox}

% ─────────────────────────────────────────────────────────────
\subsection{ب.۳ — سنت تحلیلی و قاره‌ای: یک گفتگوی ناتمام}
\label{subsec:analytic_continental}
% ─────────────────────────────────────────────────────────────

یکی از مهم‌ترین شکاف‌های فلسفه‌ی معاصر، شکاف بین سنت \concept{تحلیلی} و \concept{قاره‌ای} است. هر دو درباره‌ی آزادی حرف می‌زنند، اما نه به یک زبان.

\begin{table}[h!]
\centering
\caption{مقایسه‌ی رویکرد تحلیلی و قاره‌ای به مفهوم آزادی}
\label{tab:analytic_continental}
\renewcommand{\arraystretch}{1.6}

\begin{tabularx}{\textwidth}{|>{\bfseries}p{2.8cm}
  |>{\raggedright}X|>{\raggedright\arraybackslash}X|}
\hline
\rowcolor{PrimaryDeep}
\textcolor{white}{\textbf{معیار}} &
\textcolor{white}{\textbf{سنت تحلیلی}} &
\textcolor{white}{\textbf{سنت قاره‌ای}} \\
\hline

\rowcolor{BgCool}
روش & تحلیل منطقی مفاهیم، وضوح تعریف & تفسیر، پدیدارشناسی، گفتمان \\
\hline

\rowcolor{white}
پرسش اصلی & «آزادی دقیقاً یعنی چه؟» & «آزادی چه تجربه‌ای است؟» \\
\hline

\rowcolor{BgCool}
متفکران شاخص & برلین، رالز، پتیت، سن، مکالوم & هایدگر، سارتر، فوکو، باتلر \\
\hline

\rowcolor{white}
نسبت با دولت & مستقیم و سیاسی & غالباً انتقادی–غیرمستقیم \\
\hline

\rowcolor{BgCool}
نگاه به ساختار & غالباً نهادگرا & غالباً ساختارانتقادی \\
\hline

\rowcolor{white}
نقطه‌ی ضعف & انتزاع از واقعیت زنده‌ی قدرت & ابهام در توصیه‌های سیاسی \\
\hline

\rowcolor{BgCool}
نقطه‌ی قوت & دقت مفهومی، قابل آزمون & حساسیت به قدرت پنهان \\
\hline

\end{tabularx}
\end{table}

در دهه‌های اخیر تلاش‌هایی برای گفتگوی این دو سنت صورت گرفته. \thinker{هابرماس} با «نظریه‌ی کنش ارتباطی» (\lr{Communicative Action}) تلاش کرد دقت سنت تحلیلی را با عمق سنت قاره‌ای ترکیب کند: آزادی نه در ذهن فردی بلکه در فضای گفتگوی غیرتحریف‌آمیز (ایده‌آل گفتاری) تحقق می‌یابد.\footnote{\lr{Habermas, J. (1984). \textit{The Theory of Communicative Action}. Boston: Beacon Press.}} % FIXED: \lr in footnote

\newpage

% ─────────────────────────────────────────────────────────────
\subsection{ب.۴ — آزادی در مکاتب راست، چپ و میانه: تحلیل تقابلی}
\label{subsec:left_right_center}
% ─────────────────────────────────────────────────────────────

\begin{tcolorbox}[defbox, title={چارچوب تحلیل}]
یکی از کلیدی‌ترین محورهای تقابل در نظریه‌ی آزادی، جایگاه آن در طیف چپ–میانه–راست است. این تقابل صرفاً سیاسی نیست؛ بلکه بر سر تعریف بنیادین آزادی، نسبت آن با برابری، و نقش دولت است.
\end{tcolorbox}

\vspace{0.4cm}

\begin{table}[h!]
\centering
\caption{آزادی از منظر طیف‌های فکری سیاسی}
\label{tab:left_right}
\renewcommand{\arraystretch}{1.9}
\setlength{\tabcolsep}{4pt}

\begin{tabularx}{\textwidth}{|>{\bfseries\color{PrimaryMid}}p{2.2cm}
  |>{\raggedright}X|>{\raggedright}X|>{\raggedright}X
  |>{\raggedright\arraybackslash}X|}

\hline
\rowcolor{PrimaryDeep}
\textcolor{white}{\textbf{معیار}} &
\textcolor{white}{\textbf{راست لیبرتاریان}} &
\textcolor{white}{\textbf{لیبرال میانه}} &
\textcolor{white}{\textbf{سوسیال‌دموکراسی}} &
\textcolor{white}{\textbf{چپ رادیکال}} \\
\hline

\rowcolor{BgCool}
تعریف آزادی &
رهایی از مداخله‌ی دولت &
خودفرمانی در قوانین عادلانه &
برابری شرایط واقعی &
رهایی از سرمایه‌داری \\
\hline

\rowcolor{white}
نوع آزادی &
منفی + بازار آزاد &
منفی + مثبت محدود &
مثبت + برابری &
رهایی (Emancipation) \\
\hline

\rowcolor{BgCool}
نقش دولت &
حداقلی: پاسبان حقوق &
تنظیم‌گر: برطرف‌کننده‌ی موانع &
رفاهی: توزیع‌کننده‌ی قابلیت‌ها &
ابزار طبقه یا کمون خودگردان \\
\hline

\rowcolor{white}
رابطه با برابری &
برابری در برابر قانون (نه نتیجه) &
برابری فرصت‌ها &
برابری توانایی‌های پایه &
برابری ساختاری اقتصادی \\
\hline

\rowcolor{BgCool}
دشمن اصلی آزادی &
دولت بزرگ + مالیات &
تبعیض + استبداد &
فقر + نابرابری + پاترنالیسم &
سرمایه‌داری + امپریالیسم \\
\hline

\rowcolor{white}
نظریه‌پرداز شاخص &
هایک، نوزیک، فریدمن &
رالز، دوورکین، برلین &
گیدنز، دنی رودریک، سن &
مارکس، گرامشی، فانون \\
\hline

\rowcolor{BgCool}
ریشه‌ی انسان‌شناختی &
اتم خودبسنده، عاقل &
شخص اخلاقی با حقوق ذاتی &
موجود اجتماعی نیازمند &
موجود طبقاتی–رهایی‌طلب \\
\hline

\rowcolor{white}
نقطه‌ی کور &
نابرابری‌های قدرت اقتصادی &
بی‌توجهی به ساختارهای پنهان &
احتمال پاترنالیسم دولتی &
خطر اقتدارگرایی انقلابی \\
\hline

\end{tabularx}
\end{table}

\vspace{0.6cm}

\begin{tcolorbox}[title={کانون اصلی دعوا: آیا بازار آزاد، آزادی را گسترش می‌دهد یا محدود می‌کند؟},
  colback=AccentGold!8, colframe=AccentGold!60!black, fonttitle=\bfseries]

این سؤال بنیادی‌ترین اختلاف راست و چپ مدرن است:

\textbf{موضع راست/لیبرتاریان (هایک، فریدمن):} بازار آزاد تنها مکانیسمی است که بدون اقتدار مرکزی، هماهنگی اجتماعی را ممکن می‌کند. مداخله‌ی دولت در اقتصاد، صرف‌نظر از نیت، به سرف‌داری ختم می‌شود.\footnote{\lr{Hayek, F. (1944). \textit{The Road to Serfdom}. Chicago: University of Chicago Press.}} % FIXED: \lr in footnote

\textbf{موضع چپ/سوسیال‌دموکرات (پولانی، رالز):} قدرت بازار، شکلی از سلطه‌ی نامرئی است. کارگری که تنها گزینه‌اش پذیرش شرایط نامنصفانه‌ی کارفرماست، «آزادانه» قرارداد نمی‌بندد—او در موقعیتی شبه‌اجباری قرار دارد.\footnote{\lr{Polanyi, K. (1944). \textit{The Great Transformation}. Boston: Beacon Press, p. 257.}} % FIXED: \lr in footnote

\textbf{موضع جمهوری‌خواهانه (پتیت):} هر دو می‌توانند اشتباه کنند. آزادی نه رهایی از دولت است، نه رهایی از بازار، بلکه رهایی از \textit{سلطه‌ی دلبخواهانه}—خواه منشأ آن عمومی باشد یا خصوصی.
\end{tcolorbox}

\vspace{0.5cm}

% نمودار TikZ: طیف آزادی
\begin{center}
\begin{tikzpicture}[
  font=\small,
  every node/.style={align=center},
  arr/.style={-{Stealth[length=4pt]}, thick}
]

% خط اصلی
\draw[line width=2pt, PrimaryDeep!40] (-7,0) -- (7,0);

% علامت‌گذاری
\foreach \x/\lbl/\col in {
  -6/{\rl{لیبرتاریانیسم}\\% FIXED: گره چندخطی
  \rl{بازار آزاد}%
  }/AccentRed,
  -3/{\rl{محافظه‌کاری}\\%
  \rl{لیبرال}%
  }/AccentAmber,
   0/{\rl{لیبرالیسم}\\%
  \rl{میانه}%
  }/AccentTeal,
   3/{\rl{سوسیال}\\%
  \rl{دموکراسی}%
  }/PrimaryLight,
   6/{\rl{سوسیالیسم}\\%
  \rl{مارکسیستی}%
  }/PrimaryDeep%
}{
  \draw[\col, thick] (\x, -0.3) -- (\x, 0.3);
  \node[below=0.35cm, \col, font=\scriptsize\bfseries] at (\x, 0) {\lbl};
}

% برچسب آزادی
\node[above=0.15cm, PrimaryDeep, font=\scriptsize] at (-5, 0) {\rl{آزادی منفی}}; % FIXED: \rl
\node[above=0.15cm, AccentTeal, font=\scriptsize] at (0, 0) {\rl{توازن}}; % FIXED: \rl
\node[above=0.15cm, PrimaryMid, font=\scriptsize] at (5, 0) {\rl{آزادی مثبت}}; % FIXED: \rl

% فلش‌های راهنما
\draw[arr, AccentRed!70] (-7, 0.8) -- (-4, 0.8)
  node[midway, above, font=\scriptsize] {\rl{تأکید بر آزادی منفی}}; % FIXED: \rl
\draw[arr, PrimaryDeep!70] (4, 0.8) -- (7, 0.8)
  node[midway, above, font=\scriptsize] {\rl{تأکید بر آزادی مثبت}}; % FIXED: \rl

\end{tikzpicture}
\end{center}

% ─────────────────────────────────────────────────────────────
\subsection{ب.۵ — آزادی در سنت پسااستعماری و فمینیستی}
\label{subsec:postcolonial_feminist}
% ─────────────────────────────────────────────────────────────

\begin{tcolorbox}[defbox, title={بازتعریف آزادی از بیرون از مرکز}]
سنت‌های پسااستعماری و فمینیستی، نه تنها محتوای تعریف آزادی، بلکه \textit{چه کسی} و \textit{از کجا} آن را تعریف می‌کند را به چالش کشیده‌اند. این سنت‌ها نشان دادند که «آزادی جهانی» ادعا‌شده در فلسفه‌ی لیبرال، اغلب آزادی مرد–سفیدپوست–اروپایی بوده است.
\end{tcolorbox}

\paragraph{فانون و آزادی استعمارزده:}
\thinker{فرانتس فانون} در \lr{\textit{The Wretched of the Earth}} (۱۹۶۱) استدلال کرد که استعمار نه فقط آزادی سیاسی، بلکه سوبژکتیویته‌ی انسانی ملت‌های مستعمره را سلب کرده است. آزادی مستلزم «رهایی روانی»—شکست هویت ساخته‌شده توسط استعمارگر—پیش از هر آزادی سیاسی است.\footnote{\lr{Fanon, F. (1961/2004). \textit{The Wretched of the Earth}. New York: Grove Press, p. 2.}} % FIXED: \lr in footnote

\paragraph{اسپیواک و سکوت زیردست:}
\thinker{گایاتری اسپیواک} در مقاله‌ی «آیا زیردست می‌تواند سخن بگوید؟» (۱۹۸۸) نشان داد که ساختارهای گفتمانی چگونه صدای گروه‌های حاشیه‌ای را از پیش ناشنیدنی می‌کنند. آزادی بیان بدون ساختارهای شنیدن، ناقص است.\footnote{\lr{Spivak, G. C. (1988). Can the Subaltern Speak? In C. Nelson \& L. Grossberg (Eds.), \textit{Marxism and the Interpretation of Culture}. Urbana: University of Illinois Press, pp. 271–313.}} % FIXED: \lr in footnote

\paragraph{فمینیسم و آزادی بدن:}
\thinker{سیمون دو بووار} در \lr{\textit{The Second Sex}} (۱۹۴۹) نشان داد که «آزادی» مردانه بر مبنای تعریف زن به عنوان «دیگری» (\lr{the Other}) ساخته شده است. آزادی واقعی مستلزم از دست دادن این ساختار دوگانه‌ی سوژه–ابژه است.\footnote{\lr{de Beauvoir, S. (1949/2011). \textit{The Second Sex}. New York: Vintage Books, p. 6.}} % FIXED: \lr in footnote and text

\thinker{جودیت باتلر} این انتقاد را به حوزه‌ی هویت جنسی گسترش داد: جنسیت یک عملکرد تکراری (\lr{performativity}) است، نه یک ذات ثابت. آزادی به معنای برهم زدن این اجرا‌های تکراری است.\footnote{\lr{Butler, J. (1990). \textit{Gender Trouble: Feminism and the Subversion of Identity}. New York: Routledge, p. 25.}} % FIXED: \lr in footnote and text

\vspace{0.4cm}

\begin{table}[h!]
\centering
\caption{گسترش انتقادی مفهوم آزادی: سنت‌های پسااستعماری و فمینیستی}
\label{tab:postcolonial}
\renewcommand{\arraystretch}{1.7}

\begin{tabularx}{\textwidth}{|>{\bfseries\color{AccentTeal}}p{2.4cm}
  |>{\raggedright}X|>{\raggedright}X|>{\raggedright\arraybackslash}X|}

\hline
\rowcolor{AccentTeal!80}
\textcolor{white}{\textbf{متفکر}} &
\textcolor{white}{\textbf{انتقاد اصلی از آزادی لیبرال}} &
\textcolor{white}{\textbf{تعریف جایگزین}} &
\textcolor{white}{\textbf{برنامه‌ی سیاسی}} \\
\hline

\rowcolor{AccentTeal!10}
فانون &
فراموشی آسیب روانی استعمار &
رهایی روانی–هویتی &
انقلاب ضداستعماری \\
\hline

\rowcolor{white}
اسپیواک &
گفتمان که صداها را می‌بلعد &
آزادی گفتمانی + نمایندگی &
بازسازی ساختارهای گوش‌دهی \\
\hline

\rowcolor{AccentTeal!10}
دو بووار &
آزادی مردانه بر قربانی زن &
آزادی متقابل (Reciprocal) &
برابری جنسیتی + استقلال \\
\hline

\rowcolor{white}
باتلر &
هویت ثابت ابزار قدرت است &
آزادی اجرایی–تکرارشکنانه &
سیاست کوئیر + اغتشاش جنسیتی \\
\hline

\rowcolor{AccentTeal!10}
بِل هوکس &
آزادی سفیدِ فمینیسم موج دوم &
آزادی تقاطعی (Intersectional) &
فمینیسم ضدنژادپرست \\
\hline

\end{tabularx}
\end{table}

% ─────────────────────────────────────────────────────────────
\subsection{ب.۶ — آزادی در عصر دیجیتال: چالش‌های معاصر}
\label{subsec:digital_freedom}
% ─────────────────────────────────────────────────────────────

\begin{tcolorbox}[defbox, title={آزادی در قرن بیست‌ویکم: سه چالش جدید}]
انقلاب دیجیتال، هوش مصنوعی، و نظارت همه‌گیر، سه چالش بی‌سابقه برای هر نظریه‌ی آزادی ایجاد کرده‌اند که با چارچوب‌های کلاسیک به‌سختی قابل تحلیل هستند.
\end{tcolorbox}

\paragraph{یک: معماری پنهان انتخاب (Nudge Architecture):}
\thinker{تالر} و \thinker{سانستین} در \lr{\textit{Nudge}} (۲۰۰۸) نشان دادند که چطور طراحی محیط انتخاب می‌تواند رفتار را بدون اجبار رسمی هدایت کند. این «پدرسالاری آزادی‌پرور» (\lr{Libertarian Paternalism}) چالشی برای تعریف کلاسیک آزادی ایجاد می‌کند: آیا وقتی معماری پنهانی انتخاب ما را شکل می‌دهد هنوز «آزادانه» تصمیم می‌گیریم؟\footnote{\lr{Thaler, R. \& Sunstein, C. (2008). \textit{Nudge: Improving Decisions about Health, Wealth and Happiness}. New Haven: Yale University Press.}} % FIXED: \lr in footnote and text

\paragraph{دو: سرمایه‌داری نظارتی (Surveillance Capitalism):}
\thinker{شوشانا زوبوف} در \lr{\textit{The Age of Surveillance Capitalism}} (۲۰۱۹) استدلال کرد که اقتصاد داده، یک شکل کاملاً جدید از قدرت ایجاد کرده که تجربه‌ی انسانی را به عنوان مواد خام استخراج و به رفتار قابل پیش‌بینی تبدیل می‌کند. این نه سلطه‌ی دولت است نه بازار سنتی، بلکه چیزی است که مفاهیم آزادی کلاسیک برای توصیف آن ساخته نشده‌اند.\footnote{\lr{Zuboff, S. (2019). \textit{The Age of Surveillance Capitalism}. New York: PublicAffairs, p. 8.}} % FIXED: \lr in footnote

\paragraph{سه: آزادی الگوریتمی و تبعیض پنهان:}
تبعیض الگوریتمی نشان داده که سیستم‌های به‌ظاهر بی‌طرف می‌توانند شکل‌های تاریخی نابرابری را بازتولید کنند. برای رویکرد قابلیتی سن، این به معنای محرومیت قابلیتی غیرمرئی است که در معماری زیرساخت‌های دیجیتال تعبیه شده.

% ─────────────────────────────────────────────────────────────
\subsection{ب.۷ — جمع‌بندی: نقشه‌ی تنش‌های حل‌نشده}
\label{subsec:summary_tensions}
% ─────────────────────────────────────────────────────────────

\begin{center}
\begin{tikzpicture}[
  font=\small,
  node distance=1.4cm and 2.2cm,
  box/.style={
    rectangle, rounded corners=6pt,
    draw=PrimaryDeep, thick,
    fill=PrimaryDeep!8,
    text width=3.2cm,
    minimum height=1.1cm,
    align=center,
    font=\small\bfseries
  },
  tensionbox/.style={
    rectangle, rounded corners=4pt,
    draw=AccentRed!70, thick,
    fill=AccentRed!6,
    text width=3cm,
    minimum height=0.9cm,
    align=center,
    font=\scriptsize
  },
  arr/.style={
    -{Stealth[length=5pt]},
    AccentAmber!80!black, thick
  },
  dblarr/.style={
    {Stealth[length=4pt]}-{Stealth[length=4pt]},
    PrimaryMid!60, thick, dashed
  }
]

% گره‌های اصلی
\node[box, fill=AccentGold!15, draw=AccentGold!80!black] (freedom) at (0,0)
  {\rl{آزادی}\\% FIXED: گره چندخطی
  \rl{مفهوم مرکزی}%
  };

\node[box] (negative) at (-4.5, 1.8) {\rl{آزادی منفی}\\%
  \rl{(برلین، هایک)}%
  };
\node[box] (positive) at ( 4.5, 1.8) {\rl{آزادی مثبت}\\%
  \rl{(روسو، هگل)}%
  };
\node[box] (republican) at (-4.5,-1.8) {\rl{غیر‌سلطه}\\%
  \rl{(پتیت، آرنت)}%
  };
\node[box] (capability) at ( 4.5,-1.8) {\rl{قابلیت}\\%
  \rl{(سن، نوسبام)}%
  };
\node[box] (emancipation) at (0, 3.5) {\rl{رهایی}\\%
  \rl{(مارکس، فانون)}%
  };
\node[box] (soteriological) at (0,-3.5) {\rl{آزادی دینی}\\%
  \rl{(اسلام، مسیحیت)}%
  };

% تنش‌ها
\node[tensionbox] (t1) at (-2.5, 0.4) {\rl{فرد}\\% FIXED: گره چندخطی
  \rl{در برابر جمع}%
  };
\node[tensionbox] (t2) at ( 2.5, 0.4) {\rl{صوری}\\%
  \rl{در برابر واقعی}%
  };
\node[tensionbox] (t3) at (0, 1.5)    {\rl{اقتدار}\\%
  \rl{در برابر آزادی}%
  };
\node[tensionbox] (t4) at (-2.5,-1.2) {\rl{هویت}\\%
  \rl{در برابر انتخاب}%
  };
\node[tensionbox] (t5) at ( 2.5,-1.2) {\rl{مسئولیت}\\%
  \rl{در برابر ساختار}%
  };

% اتصالات
\draw[dblarr] (freedom) -- (negative);
\draw[dblarr] (freedom) -- (positive);
\draw[dblarr] (freedom) -- (republican);
\draw[dblarr] (freedom) -- (capability);
\draw[dblarr] (freedom) -- (emancipation);
\draw[dblarr] (freedom) -- (soteriological);

% تیترهای تنش
\node[above=0.05cm, AccentRed!80, font=\tiny\bfseries] at (-2.5, 0.9) {\rl{تنش ۱}}; % FIXED: \rl
\node[above=0.05cm, AccentRed!80, font=\tiny\bfseries] at ( 2.5, 0.9) {\rl{تنش ۲}}; % FIXED: \rl
\node[above=0.05cm, AccentRed!80, font=\tiny\bfseries] at (0.4, 1.9)  {\rl{تنش ۳}}; % FIXED: \rl
\node[above=0.05cm, AccentRed!80, font=\tiny\bfseries] at (-2.5,-0.8) {\rl{تنش ۴}}; % FIXED: \rl
\node[above=0.05cm, AccentRed!80, font=\tiny\bfseries] at ( 2.5,-0.8) {\rl{تنش ۵}}; % FIXED: \rl

\end{tikzpicture}
\end{center}

\vspace{0.5cm}

\begin{tcolorbox}[
  title={نتیجه‌گیری پیوست ب: آزادی—مفهومی که هیچ‌گاه بسته نمی‌شود},
  colback=PrimaryDeep!5, colframe=PrimaryDeep!60, fonttitle=\bfseries,
  breakable
]

فلسفه‌ی آزادی در وضعیتی است که \thinker{گالی} آن را «مفهوم ذاتاً مورد مناقشه» (\lr{Essentially Contested Concept}) نامید: مفهومی که ارزیابی‌های رقیب نه از سوء‌تفاهم، بلکه از تعهدات ارزشی متفاوت ناشی می‌شوند.\footnote{\lr{Gallie, W. B. (1956). Essentially Contested Concepts. \textit{Proceedings of the Aristotelian Society}, 56, 167–198.}} % FIXED: \lr in footnote

این بدان معنا نیست که «هر تعریفی درست است.» بلکه به این معناست که:

\begin{enumerate}[leftmargin=2em, itemsep=4pt]
  \item هر نظریه‌ی آزادی باید درباره‌ی \textbf{انسان‌شناسی} خود صادق باشد.
  \item هر نظریه باید \textbf{تنش‌های درونی} خود را بشناسد و آدرس دهد.
  \item هیچ نظریه‌ای نمی‌تواند ادعای \textbf{جامعیت بدون قربانی‌کردن برخی ارزش‌ها} داشته باشد.
  \item کار نظریه‌پرداز آزادی نه حل نهایی، بلکه \textbf{شفاف‌سازی معاوضه‌ها} (trade-offs) است.
\end{enumerate}

در سنت تحلیلی، این صداقت مفهومی خود نوعی آزادی فکری است—آزادی از توهم نظام‌های بسته.

\end{tcolorbox}

% ─────────────────────────────────────────────────────────────
% پایان پیوست ب
% ─────────────────────────────────────────────────────────────
